\chapter{검증과 정정의 기록}

\section{우연을 지우는 도구들}
\begin{table}[H]\centering\small
\begin{tabularx}{\linewidth}{@{}p{38mm}LL@{}}\toprule
의심 & 도구 & 결과\\\midrule
운이 좋았던 것 아닌가 & 시드 3--10개, 발행월 단위 부트스트랩(1만 번, 6개월 이동 블록) & 구간이 0에서 멀다\\
평가 코드가 가짜 성과를 만드나 & 정답을 무작위로 섞은 가짜 모델, 무작위 신호 롱숏 & 0 근처(정상)\\
미래 정보가 새지 않았나 & 컬럼 시점 점검, 시점 학습, 상각 지연 5개월, 약한 누수 후보(투자자 펀딩액) 제거 & 결과 동일\\
모델·\allowbreak{}전처리 덕인가 & 모델 17종, 결측 처리 9종, 피처 6종 & 거의 같음\\
채점 표본이 편향됐나 & 다 끝난 대출을 일부러 일찍 자르는 시뮬레이션 & 편향 확인, test 해석 규칙 고정\\
결과를 보고 기준을 바꾸지 않았나 & 판정 기준을 스크립트 맨 위에 먼저 적음, test 규칙 git 태그 & 실패도 실패로 기록\\
여러 번 비교한 탓 아닌가 & Holm 보정, 시도한 조합 수 공개(80개 이상) & 핵심 결과는 보정 후에도 유의\\\bottomrule
\end{tabularx}
\end{table}

\section{처음 판단이 틀렸던 것}
\begin{table}[H]\centering\small
\begin{tabularx}{\linewidth}{@{}LLp{22mm}@{}}\toprule
처음 판단 & 실제 & 어떻게 알았나\\\midrule
부도확률 순위는 무작위와 비슷하다 & 시드 여러 개로 보면 유의하게 양수 & 14단계\\
선별 곡선은 새로운 방법이다 & 수익 기반 점수(2016), EMP(2014)가 이미 있음 & 문헌 검색\\
로지스틱이 무너진 건 결측 처리 탓 & 학습 때 거의 없던 변수의 표준화 폭발 & 22단계\\
CatBoost 붕괴는 새 컬럼이 조기 종결의 대리 변수라서 & 부도율 17.2\% vs 17.6\%로 같아 기각, 해법(가용성 규칙)만 확인 & 24b단계\\
$p$의 크기가 맞아야 돈을 번다 & 극단적으로 틀릴 때만 문제 & 25단계\\
2017년에는 부도확률 순위가 더 낫다 & 끝난 대출만 본 종결 편향 & 26단계\\
\textbf{Sharpe는 승인비율에 끌려가 주 지표로 부적합} & 대출 건별 Sharpe에만 맞음. 포트폴리오 Sharpe에서는 보수적 선별이 우위(팀원 지적) & 30단계\\
\textbf{LC 등급끼리의 가격은 공정했다} & A--D 평균 수익에만 맞음. 36개월 E--G는 금리가 위험을 못 메움 & 33단계\\
고위험 등급은 손해다 & 36개월에만 맞음. 60개월은 반대 & 34단계\\
단순 규칙이 수익·\allowbreak{}Sharpe 모두 1위 & 동점 버그(실제 승인 39.5\%). 고치면 Sharpe 우위만 & 34b단계\\
\textbf{기대수익 순은 무작위보다 Sharpe가 낮다} & 발행연도 단위 컷이 만든 값. 발행월 단위면 2.91로 무작위(2.69)보다 높다. 철회 & 외부 검증, 37단계\\
\textbf{안정성은 부도확률 낮은 순이 1위} & 허들이 국채 수준일 때만. BAA 기준이면 기대수익 순 & 외부 검증, 37--38단계\\
\textbf{1-SE 기준선 안 = 최고와 통계적으로 구분 안 됨} & 1위 하나의 SE로 그은 선. 짝지은 비교로는 10\% $-$ 80\% = 0.76(0.45--1.60)으로 구분됨. 경험적 선택 규칙으로 정정 & 팀원 검토, 39단계 재점검\\
\textbf{승인비율이 달라도 평균과 표준편차는 비례할 뿐} & 대출 구성도 바뀐다(Sharpe 3.82 \ensuremath{\to} 3.06) & 팀원 검토, 41단계 재점검\\
기대수익 순은 작은 예산에서도 무작위보다 덜 흔들린다 & 10\%에서 Sharpe 2.55 vs 2.63, 50\% 이하 구분 안 됨. 60\% 이상에서만 높음 & 42단계\\
과제가 Sharpe를 요구하니 Sharpe 기준으로 최종안을 고른다 & 공지 재확인 결과 필수 아님. 최종안을 기대수익 순(v1 H2)으로 변경 & 공지 재확인, v2.3\\\bottomrule
\end{tabularx}
\caption{정정 기록(\ft{docs/research_log.md} 3절의 요약).}
\end{table}

\begin{insight}[정정에서 배운 것]
정정의 절반 이상이 ``숫자는 맞는데 범위를 너무 넓게 말한'' 경우였다. Sharpe 주장은 건별 정의에만, 등급 가격 주장은 A--D에만, 고위험 손해 주장은 36개월에만 맞았다. 그래서 이제 모든 주장에 \textbf{어느 표본, 어느 정의에서 맞는지}를 붙인다.
\end{insight}

\begin{warnbox}[스스로 검증할 때 왜 놓쳤나]
34단계에서 ``근본 가정 재점검''을 했지만 위의 마지막 두 정정은 외부 검증자가 찾았다. (1) 30--36단계가 같은 승인 함수(발행연도 단위)를 재사용했고, 다시 확인할 때도 그 함수로 돌려 함수 안의 선택이 의심받지 않았다. (2) 본문의 조건 설명(``같은 달'')을 코드가 아니라 의도대로 썼다. (3) 점검 목록에 비교 기준(허들)을 넣지 않았다. 외부 검증자는 본문만 보고 처음부터 따로 구현해서 차이를 찾았다. 그래서 이제 핵심 숫자는 독립 구현으로 한 번 더 확인하고, 문서의 조건은 코드와 대조한다.
\end{warnbox}

\section{데이터 엿보기의 위험}
같은 2013--16년 구간을 1단계부터 36단계까지 계속 다시 썼다. 판정 기준을 미리 정했지만, 수십 번 들여다본 구간에서 좋아 보이는 것은 그 구간에 맞춰졌을 수 있다. 특히 단순 규칙(7장)은 결과를 본 뒤 고른 것이다. 이 위험을 줄이는 유일한 방법은 한 번도 보지 않은 공식 test로 미리 정한 규칙을 한 번만 돌리는 것이다(10장).
